You are given a bit string of length $n$. Your task is to calculate for each $k$ between $0  \ldots n$ the number of non-empty substrings that contain exactly $k$ ones.
For example, if the string is 101, there are:

\textbullet\ 1 substring that contains 0 ones: 0
\textbullet\ 4 substrings that contain 1 one: 01, 1, 1, 10
\textbullet\ 1 substring that contains 2 ones: 101
\textbullet\ 0 substrings that contain 3 ones

\vspace{0.3cm}\noindent\textbf{Input}

The only input line contains a binary string of length $n$.

\vspace{0.3cm}\noindent\textbf{Output}

Print $n+1$ values as specified above.

\vspace{0.3cm}\noindent\textbf{Constraints}

\textbullet\ $1 \le n \le 2 \cdot 10^5$